Queremos sumar los números binarios $010\ (2)$ y $110\ (3)$.
Necesitamos cuatro ancillas y por lo tanto el circuito tiene diez cúbits.

De acuerdo con la convención postfija, el circuito se lee de izquierda a derecha:
\begin{equation*}
  C = \stategate{X}{1}{3}\stategate{X}{5}{7}\stategate{X}{2}{3}\stategate{X}{6}{7}\stategate{X}{1,2}{4}\stategate{X}{5,6}{8}\stategate{X}{4}{9}\stategate{X}{7}{9}\stategate{X}{4,7}{10}\stategate{X}{8}{10}
\end{equation*}

y opera sobre el ket $\ket{a}=\ket{010}\otimes\ket{0}\otimes\ket{110}\otimes\ket{0}\otimes\ket{00}=\ket{0100110000}$ por pasos:

\begin{align*}
  \ket{a}C^1    = \ket{0100110000}\stategate{X}{1}{3}    & = \ket{0100110000} \\
  \ket{a}C^2    = \ket{0100110000}\stategate{X}{5}{7}    & = \ket{0100111000} \\
  \ket{a}C^3    = \ket{0100111000}\stategate{X}{2}{3}    & = \ket{0110111000} \\
  \ket{a}C^4    = \ket{0110111000}\stategate{X}{6}{7}    & = \ket{0110110000} \\
  \ket{a}C^5    = \ket{0110110000}\stategate{X}{1,2}{4}  & = \ket{0110110000} \\
  \ket{a}C^6    = \ket{0110110000}\stategate{X}{5,6}{8}  & = \ket{0110110100} \\
  \ket{a}C^7    = \ket{0110110100}\stategate{X}{4}{9}    & = \ket{0110110100} \\
  \ket{a}C^8    = \ket{0110110100}\stategate{X}{7}{9}    & = \ket{0110110100} \\
  \ket{a}C^9    = \ket{0110110100}\stategate{X}{4,7}{10} & = \ket{0110110100} \\
  \ket{a}C^{10} = \ket{0110110100}\stategate{X}{8}{10}   & = \ket{0110110101} \\
\end{align*}

Y ahora medimos sobre los cúbits en las posiciones 3, 9 y 10:
\begin{equation*}
  _{\set{3,9,10}}\rfloor\ket{0110110101}=101\ (5)
\end{equation*}
